\section{Early Playtest Findings}

Data was collected through PostHog, an open-source product analytics platform embedded in the game client. No identifying information was requested or collected. We assigned each browser session an anonymous distinct ID, and we recorded a small set of gameplay events (\texttt{game\_started}, \texttt{avatar\_selected}, \texttt{decision\_made}, \texttt{proworker\_stance\_chosen}, \texttt{policy\_voted}, and \texttt{ending\_reached}) along with the choice each player made at branch points. The reporting window covers May 1 through May 10, 2026.

Across that window, 40 unique players started 88 sessions. About 90\% of starters (36 players) made it past the intro and avatar selection, and about 74\% (29 players) reached at least one of the four endings. Run-throughs totaled 46, since some players finished more than once and arrived at different endings on different attempts. The distribution across the four pathways was uneven: about 52\% of completed runs ended in Augmentation, 30\% in Pro-Worker Redesign, 11\% in Selective Shrinkage, and 7\% in Job Erosion. The two outcomes that involve clear worker harm (Erosion and Shrinkage) together made up roughly 18\% of completed runs. Because the game's branch routing mixes player choices with some designed randomness in the routing logic, this distribution should not be read as a direct mapping of player intent, but the majority of playthroughs did end on a path where workers retain more agency.

Both in-game ballots produced lopsided results. The AI Job Preservation and Work Sharing Act, which appears in 2030 on every path, drew 42 votes: about 81\% in favor, about 17\% against, and one abstention (about 2\%). All seven against votes came from players whose runs ended in Job Erosion or Selective Shrinkage; players on the two pro-worker paths voted yes in every case except the single abstention. The AI Dividend Act in 2032, which only surfaces on the pro-worker path, received 14 votes, all of them yes, even though the in-game framing tells employed players they would receive a smaller share of the proceeds. The sample is small and self-selected toward people willing to spend 20 to 30 minutes on a scenario about AI and labor, but among those who did finish, support for labor protections and for redistributive AI policy was close to unanimous.
